\scanpage{37}
% Pencil archive foliation: 36.
\uuline{longueur finie}. Elle est stable par $\otimes$,
sur cette catégorie les $\operatorname{Ext}^i$ $(i>0)$ sont
nuls, et la catégorie est aussi stable
par $\operatorname{Hom}$ (cas de $\mathcal{M}(X)_{\mathrm{spec.}}$). On a
une ``catégorie tensorielle sur $\mathbb{Q}$''.

Le \deleted{\unsure{le|la}} $\lambda$-anneau défini par cette
catégorie est noté $\mathrm{M}^+(X)$ resp. $\mathrm{M}(X)$.
Noter que $\xi=\operatorname{cl}\text{ de }\unsure{\mathbb{Q}|\mathbb{Z}}_X(-1)$ est une
\unsure{droite} de $\mathrm{M}^+(X)$, et on a
\[
\mathrm{M}(X)=\mathrm{M}^+(X)_\xi=\mathrm{M}^+(X)\left[\frac1\xi\right].
\]
En fait, $\xi$ est non div. de $0$ :
$\mathrm{M}^+(X)$ comme $\mathbb{Z}$-module admet comme
base les classes d'\deleted{\uline{éléments}}\inserted{objets} \uline{simples} de
$\mathcal{M}^+(X)$, et $x\unsure{\rightsquigarrow|\mapsto}\xi.x$ dans $\mathrm{M}^+(X)$ est défini
par une application \uline{injective} sur l'ens des
objets simples. [N.B. $M$ simple \inserted{ds $\mathcal{M}^+(X)$}
$\Longleftrightarrow$ \inserted{$M$ simple dans $\mathcal{M}(X)\Longleftrightarrow$}
$M(-1)$ simple, comme il résulte du
fait que $\mathcal{M}^+(X)$ est épaisse dans $\mathcal{M}(X)$ et
$\otimes\mathbb{Z}(-1)$ une autoéquivalence.]
